\chapter{Conclusões e Trabalhos Futuros}
\label{chap:conclusoes-e-trabalhos-futuros}

O presente trabalho investigou o potencial e a viabilidade da utilização de
\gls{LLM}s no contexto da \gls{PCG} para jogos do gênero \textit{roguelike}. A
pesquisa partiu da identificação de uma lacuna significativa na literatura e nas
capacidades técnicas atuais: embora os \gls{LLM}s demonstrem proficiência
semântica e criativa, eles possuem limitações intrínsecas no raciocínio espacial
e geométrico, o que dificulta a geração direta de topologias de mapas complexos
e funcionais.

Para solucionar este problema, foi proposta e implementada uma arquitetura
híbrida baseada no desacoplamento de responsabilidades. O sistema delegou a
construção estrutural e topológica da masmorra a algoritmos de \gls{PCG}
tradicionais e determinísticos, garantindo a jogabilidade e a conectividade das
salas.  Simultaneamente, atribuiu-se ao \gls{LLM} a função de "agente criativo",
responsável pela geração semântica de ativos, incluindo a definição de temas,
narrativas, inimigos e itens. Esta abordagem permitiu contornar as deficiências
espaciais dos modelos de linguagem, aproveitando sua capacidade de interpretar
intenções textuais e enriquecer o contexto do jogo.

Um componente central para o sucesso desta implementação foi a utilização da
técnica de \gls{RAG} aliada ao uso de \textit{Structured Outputs}. A segmentação
do banco vetorial em três categorias distintas (\textit{Entities},
\textit{Items} e \textit{Environments}) provou-se uma estratégia eficaz para
mitigar alucinações visuais, impedindo, por exemplo, a seleção de texturas de
arquitetura para representar personagens. A curadoria e anotação manual deste
\textit{dataset}, composto por 489 pares de descrições e imagens, constitui uma
contribuição técnica relevante deste trabalho, estabelecendo uma base sólida
para a associação entre linguagem natural e ativos de \textit{pixel art}.

A validação do sistema, realizada através de uma abordagem quantitativa e
qualitativa, corroborou a eficácia do \textit{pipeline} proposto. A aplicação da
métrica de Reconstrução Semântica, baseada na similaridade de cosseno,
demonstrou que o sistema é capaz de preservar a intenção original do usuário ao
traduzir descrições textuais em objetos de jogo. Mais notavelmente, os testes
comparativos entre diferentes modelos revelaram que \gls{LLM}s compactos e
executados localmente, como o \texttt{openai/gpt-oss-20b}, apresentam desempenho
pareável a modelos de maior porte e complexidade como o \texttt{gpt-oss-120b} e
o \texttt{llama-4-maverick}, confirmando a viabilidade técnica de integrar esta
tecnologia em hardware de consumo, sem dependência mandatória de conexões com a
nuvem.

Em suma, este trabalho demonstra que a integração entre algoritmos clássicos de
\gls{PCG} e a moderna Inteligência Artificial Generativa não apenas é possível,
como também oferece um caminho promissor para o desenvolvimento de jogos. O
protótipo desenvolvido exemplifica como a barreira técnica da criação de
conteúdo pode ser reduzida, permitindo que desenvolvedores e jogadores utilizem
linguagem natural para moldar experiências de jogo únicas e coerentes.

\section{Limitações e Trabalhos Futuros}

Apesar dos resultados satisfatórios obtidos com a arquitetura proposta, o
sistema apresenta limitações que abrem caminhos para futuras investigações e
aprimoramentos. A restrição mais significativa reside na dependência de um banco
vetorial estático para a representação visual. O mecanismo de \gls{RAG}, embora
eficiente na recuperação, está limitado ao universo de texturas pré-existentes
no \textit{dataset}. Em cenários nos quais a descrição textual do usuário solicita um
objeto específico não catalogado (por exemplo, uma "foice"), o sistema seleciona
a aproximação semântica mais próxima disponível (como uma "espada"), o que pode
gerar inconsistências visuais e reduzir a liberdade criativa.

Adicionalmente, o escopo atual de geração restringe-se a entidades isoladas. O
sistema instancia objetos individuais, mas ainda não possui a capacidade de
orquestrar a geração de estruturas compostas ou arranjos arquitetônicos
complexos, como monumentos, cemitérios ou acampamentos, que exigiriam um
entendimento mais refinado de composição de cena por parte do \gls{LLM}.  Diante
destes desafios, sugerem-se as seguintes direções para trabalhos futuros:

\begin{enumerate}
    \item \textbf{Integração com IA Generativa de Imagens:} A substituição ou
    complementação do banco vetorial fixo por modelos de difusão latente (como
    \textit{Stable Diffusion}) gerando texturas em tempo real. Isso eliminaria a
    limitação do acervo de ativos, permitindo que a representação visual seja
    tão dinâmica e ilimitada quanto as descrições textuais geradas pelo
    \gls{LLM}.

    \item \textbf{Geração de Estruturas Complexas:} Expandir a lógica do
    \textit{Gerador de Assets} para que o \gls{LLM} possa definir não apenas
    objetos únicos, mas conjuntos de elementos pré-fabricados
    (\textit{prefabs}). Isso permitiria povoar o mapa com micro-cenários
    semanticamente coesos, aumentando a profundidade da ambientação.

    \item \textbf{Aplicação em Novos Gêneros:} Adaptar a metodologia para outros
    gêneros de jogos, como RPGs, Plataforma ou \textit{Puzzles}. No contexto de
    \gls{RPG}s, por exemplo, o sistema poderia ser expandido para gerar
    \gls{NPC}s com árvores de diálogos complexas, missões personalizadas e
    histórias de fundo que reagem dinamicamente ao tema escolhido pelo jogador,
    explorando ainda mais a capacidade narrativa dos modelos de linguagem.
\end{enumerate}